

%\documentclass[ngerman,xcolor={table,dvipsmames},hyperref={bookmarks,colorlinks},compress,handout]{beamer}
\documentclass[ngerman,xcolor={table,dvipsmames},hyperref={bookmarks,colorlinks},compress]{beamer}
\usepackage[ngerman]{babel}	
\usepackage{textcomp}
\usepackage{listings}	
%\usepackage{../handoutWithNotes}


\usetheme{Antibes}
\usecolortheme{seahorse}	
\useinnertheme{rounded}
%\setbeamercolor*{block title alerted}{fg=black,bg=white}

%albatross | beaver | beetle |
%	crane | default | dolphin |
%	dove | fly | lily | orchid |
%	rose |seagull | seahorse |
%	sidebartab | structure |
%	whale | wolverine

%\usetheme{Berlin}	

\logo{\includegraphics[height=0.5cm]{../logo.jpg}}

%\AtBeginSection{%
%  \begin{frame}<beamer>{Gliederung}
%      \tableofcontents[currentsection, hideallsections]
%  \end{frame}
%}






\usepackage[utf8]{inputenc}
\usepackage[T1]{fontenc}
\usepackage[ngerman]{babel}
%\usepackage{rail}
%\railoptions{-t}
%\pgfpagesuselayout{2 on 1 with notes landscape}[a4paper,border shrink=5mm]
%\pgfpagesuselayout{3 on 1 with notes}[a4paper,border shrink=5mm]
\date{26. November 2013}
\author {Alexander Pietz}
\title{Die neue Sachlichkeit}


\begin{document}
\pagenumbering{roman}
\begin{frame}
\maketitle
\end{frame}

\section{Einführung}
\begin{frame}{Inhalt}
\tableofcontents
\end{frame}

\begin{frame}{Die neue Sachlichkeit}
\begin{itemize}
\item Richtung der Literatur in der Weimarer Republik
\item schlichte Klarheit, sachliche Ausdrucksweise, politischer Inhalt
\item ``alte Sachlichkeit'': Realismus
\end{itemize}
\end{frame}

\section{Geistige Voraussetzungen}
\begin{frame}{Geistige Voraussetzungen}
\begin{itemize}
\item depressive Stimmung
\item Massen- und Mediengesellschaft
\item Armut
\item Emanzipation der Frau
\end{itemize}
\end{frame}

\section{Zeitliche Einordnung}
\begin{frame}{Zeitliche Einordnung}
\begin{itemize}
\item 20er und frühe 30er Jahre
\item Beginn fraglich
\item Ende 1933
\item Bücher verbrannt, Autoren gefangen genommen
\end{itemize}
\end{frame}

\section{Merkmale der Neuen Sachlichkeit}
\begin{frame}{Inhalte und Themen}
\begin{itemize}
\item objektive Darstellung der Realität
\item kritischer Blick auf Gesellschaft
\item Beliebte Motive:
\begin{itemize}
\item soziale, politische und wirtschaftliche Wirklichkeit der Weimarer Republik 
\item Nachwirkungen des Ersten Weltkrieges
\item Inflation
\end{itemize}
\end{itemize}
\end{frame}

\begin{frame}{Sprache}
\begin{itemize}
\item kühl und distanziert
\item einfache Alltagssprache
\item Minimum an Sprache, Maximum an Bedeutung
\end{itemize}
\end{frame}

\begin{frame}{Figuren}
\begin{itemize}
\item sachlich
\item einfache Leute aus der mordernen Massengesellschaft
\item Gefühle vorhanden, aber kaum gezeigt
\end{itemize}
\end{frame}

\section{Anliegen}
\begin{frame}{Anliegen}
\begin{itemize}
\item objektive, genaue Wiedergabe der Realität
\item Leitbilder zum Bestehen in der modernen Massen- und Mediengesellschaft
\item Wiederspiegeln von Alltagssorgen
\item Bevölkerung für Demokratie begeistern
\end{itemize}
\end{frame}

\section{Literarische Gattungen}
\begin{frame}{Gebrauchslyrik}
\begin{itemize}
\item Gedichte mit Zweck:
\item Leser auf Missstände aufmerksam machen
\item sofortige Wirkung
\end{itemize}
\end{frame}

\begin{frame}{Zeitroman}
\begin{itemize}
\item umfassende Informationen über die Zeit
\item schwierige Lebenssituationen
\item Kritik an der Zeit
\end{itemize}
\end{frame}

\begin{frame}{Andere Gattungen}
\begin{itemize}
\item Reportageliteratur
\item Episches Theater
\item Kritisches Volkstheater
\end{itemize}
\end{frame}

\section{Abschluss}

\begin{frame}{Schluss und Quellen}
Vielen Dank f"ur ihre Aufmerksamkeit!\vspace{5ex} \pause\newline
  \textbf{Quellen:}\newline
uni-due.de\newline dhm.de\newline literaturwelt.com\newline zeit.de\newline wikipedia.org \newline Suchmaschine: duckduckgo.com

\end{frame}
\end{document}
